\documentclass[11pt,a4paper]{article}
\usepackage[utf8]{inputenc}
\usepackage[vietnamese]{babel}
\usepackage{amsmath,amssymb}
\usepackage{geometry}
\geometry{a4paper, total={180mm,260mm}, left=15mm, right=15mm, top=18mm, bottom=18mm}

\usepackage{tikz}
\usetikzlibrary{shapes.geometric, arrows.meta, positioning, calc}
\tikzset{arr/.style={-{Stealth[scale=0.8]}}}
\usepackage{tcolorbox}
\tcbuselibrary{skins,breakable}
\usepackage{fontawesome}
\usepackage{fancyhdr}
\usepackage{enumitem}

\pagestyle{fancy}
\fancyhf{}
\renewcommand{\headrulewidth}{0.5pt}
\renewcommand{\footrulewidth}{0.5pt}
\fancyhead[L]{\small\bfseries\color{violet!85!black}\faBook\ SỔ TAY KIẾN THỨC TOÁN 2}
\fancyhead[R]{\small\bfseries\color{violet!85!black}Lớp Toán Thầy Lĩnh -- loptoanthaylinh.vn}
\fancyfoot[C]{\small\color{gray}-- Sổ tay kiến thức: Bài 12. Bảng trừ có nhớ trong phạm vi 20 --}

\setlength{\parindent}{0pt}
\setlength{\parskip}{0pt}

\begin{document}

% --- BANNER BÀI HỌC ---
\begin{tcolorbox}[
  enhanced,
  colback=violet!8!white,
  colframe=violet!70!black,
  arc=3mm, boxrule=1.2pt,
  left=3mm, right=3mm, top=2mm, bottom=2mm,
  halign=center
]
  {\LARGE\bfseries\color{violet!90!black}\faBookmark\ BÀI 12: BẢNG TRỪ (CÓ NHỚ) TRONG PHẠM VI 20}\\[1mm]
  {\itshape\footnotesize\color{black!75} Hệ thống các bảng trừ 11, 12, 13, 14, 15, 16, 17, 18 trừ đi một số trên 1 trang A4}
\end{tcolorbox}

\vspace{-2mm}

\noindent
\begin{minipage}[t]{0.49\textwidth}\vspace{0pt}
\begin{tcolorbox}[
  enhanced,
  colback=violet!4!white,
  colframe=violet!75!black,
  arc=2mm, boxrule=0.8pt,
  left=2mm, right=2mm, top=1.8mm, bottom=1.8mm,
  title={\footnotesize\bfseries\color{white}\faKey\ 1. Bảng trừ 11, 12, 13, 14}
]
{\footnotesize
\begin{tabular}{l|l}
  \textbf{11 trừ đi một số:} & \textbf{12 trừ đi một số:} \\ \hline
  $11 - 2 = 9$; $11 - 3 = 8$ & $12 - 3 = 9$; $12 - 4 = 8$ \\
  $11 - 4 = 7$; $11 - 5 = 6$ & $12 - 5 = 7$; $12 - 6 = 6$ \\
  $11 - 6 = 5$; $11 - 7 = 4$ & $12 - 7 = 5$; $12 - 8 = 4$ \\
  $11 - 8 = 3$; $11 - 9 = 2$ & $12 - 9 = 3$ \\ \hline
  \textbf{13 trừ đi một số:} & \textbf{14 trừ đi một số:} \\ \hline
  $13 - 4 = 9$; $13 - 5 = 8$ & $14 - 5 = 9$; $14 - 6 = 8$ \\
  $13 - 6 = 7$; $13 - 7 = 6$ & $14 - 7 = 7$; $14 - 8 = 6$ \\
  $13 - 8 = 5$; $13 - 9 = 4$ & $14 - 9 = 5$
\end{tabular}
}
\end{tcolorbox}
\end{minipage}
\hfill
\begin{minipage}[t]{0.49\textwidth}\vspace{0pt}
\begin{tcolorbox}[
  enhanced,
  colback=teal!4!white,
  colframe=teal!75!black,
  arc=2mm, boxrule=0.8pt,
  left=2mm, right=2mm, top=1.8mm, bottom=1.8mm,
  title={\footnotesize\bfseries\color{white}\faSliders\ 2. Bảng trừ 15, 16, 17, 18}
]
{\footnotesize
\begin{tabular}{l|l}
  \textbf{15 trừ đi một số:} & \textbf{16 trừ đi một số:} \\ \hline
  $15 - 6 = 9$; $15 - 7 = 8$ & $16 - 7 = 9$; $16 - 8 = 8$ \\
  $15 - 8 = 7$; $15 - 9 = 6$ & $16 - 9 = 7$ \\ \hline
  \textbf{17 trừ đi một số:} & \textbf{18 trừ đi một số:} \\ \hline
  $17 - 8 = 9$; $17 - 9 = 8$ & $18 - 9 = 9$
\end{tabular}

\vspace{1.5mm}
\textbf{Quan hệ giữa phép cộng và phép trừ:}\\
Từ phép cộng $a + b = c$, ta suy ra:
$$c - a = b \quad\text{và}\quad c - b = a$$
\textit{Ví dụ:} $8 + 7 = 15 \implies 15 - 8 = 7$ và $15 - 7 = 8$.
}
\end{tcolorbox}
\end{minipage}

\vspace{-1.5mm}
\begin{tcolorbox}[
  enhanced,
  colback=orange!4!white,
  colframe=orange!80!black,
  arc=2mm, boxrule=0.8pt,
  left=2mm, right=2mm, top=1.5mm, bottom=1.5mm,
  title={\footnotesize\bfseries\color{white}\faSitemap\ 3. Sơ đồ tư duy: Tam giác liên kết Phép cộng và Phép trừ}
]
\centering
\begin{tikzpicture}[x=1.5cm, y=1cm]
  \node[circle, draw=violet!80!black, fill=violet!20, font=\footnotesize\bfseries, inner sep=4pt] (tong) at (0,1) {$15$};
  \node[rectangle, rounded corners=2pt, draw=teal!80!black, fill=teal!15, font=\scriptsize\bfseries] (sh1) at (-1.5,-0.3) {Số hạng: $8$};
  \node[rectangle, rounded corners=2pt, draw=teal!80!black, fill=teal!15, font=\scriptsize\bfseries] (sh2) at (1.5,-0.3) {Số hạng: $7$};

  \draw[<->, thick, orange!80!black] (sh1) -- node[midway, below, font=\tiny\bfseries] {$8 + 7 = 15$} (sh2);
  \draw[->, thick, violet!80!black] (tong) -- node[left, font=\tiny\bfseries] {$15 - 8 = 7$} (sh1);
  \draw[->, thick, violet!80!black] (tong) -- node[right, font=\tiny\bfseries] {$15 - 7 = 8$} (sh2);
\end{tikzpicture}
\end{tcolorbox}

\vspace{-1.5mm}
\noindent
\begin{minipage}[t]{0.49\textwidth}\vspace{0pt}
\begin{tcolorbox}[
  enhanced,
  colback=green!4!white,
  colframe=green!60!black,
  arc=2mm, boxrule=0.8pt,
  left=2mm, right=2mm, top=1.8mm, bottom=1.8mm,
  title={\footnotesize\bfseries\color{white}\faCheckCircle\ 4. Ví dụ thực hành chuẩn mực}
]
{\footnotesize
\textbf{1. Tính nhẩm:}\\
$14 - 8 = \mathbf{6}$; $16 - 8 = \mathbf{8}$; $17 - 9 = \mathbf{8}$.\\
$13 - 7 = \mathbf{6}$; $12 - 6 = \mathbf{6}$; $18 - 9 = \mathbf{9}$.
\textbf{2. Tìm số còn thiếu:}\\
\begin{itemize}\setlength{\itemsep}{0.5pt}\setlength{\leftskip}{-2.5mm}
  \item $14 - \mathbf{5} = 9$ (vì $14 - 9 = 5$).
  \item $11 - \mathbf{6} = 5$ (vì $11 - 5 = 6$).
  \item $\mathbf{16} - 7 = 9$ (vì $9 + 7 = 16$).
\end{itemize}
}
\end{tcolorbox}
\end{minipage}
\hfill
\begin{minipage}[t]{0.49\textwidth}\vspace{0pt}
\begin{tcolorbox}[
  enhanced,
  colback=red!4!white,
  colframe=red!75!black,
  arc=2mm, boxrule=0.8pt,
  left=2mm, right=2mm, top=1.8mm, bottom=1.8mm,
  title={\footnotesize\bfseries\color{white}\faExclamationTriangle\ 5. Mẹo tránh bẫy trắc nghiệm}
]
{\footnotesize
\begin{itemize}\setlength{\itemsep}{1pt}\setlength{\leftskip}{-2.5mm}
  \item \textbf{Mẹo tìm số bị trừ:}\\
  Muốn tìm Số bị trừ, ta lấy Hiệu \textbf{cộng} với Số trừ.\\
  $\dots - 8 = 5 \implies 5 + 8 = 13$.
  \item \textbf{Mẹo tìm số trừ:}\\
  Muốn tìm Số trừ, ta lấy Số bị trừ \textbf{trừ} đi Hiệu.\\
  $15 - \dots = 9 \implies 15 - 9 = 6$.
  \item Thuộc làu các cặp đôi bằng nhau: $12-6=6, 14-7=7, 16-8=8, 18-9=9$.
\end{itemize}
}
\end{tcolorbox}
\end{minipage}


\end{document}
